\section{The MoRE Architecture}
\label{sec:method}

\subsection{Setting and notation}

The task is autoregressive language modelling on packed token sequences. Write
$E$ for the number of experts, $T$ for the maximum recursion depth, $d$ for the
model width, and $\mathbf{h}^{(0)}_s \in \mathbb{R}^{d}$ for the post-embedding
state of position $s$. Unlike a synthetic benchmark, natural language supplies
\emph{no per-token oracle}---neither an expert index nor a complexity label---so
every structural claim below must rest on permutation-invariant statistics or on
comparisons against a controlled alternative. This is why
\S\ref{sec:halting} is built around a forced-depth sweep rather than around
agreement with ground truth.

\subsection{Routing}

Routing happens once per token. A linear router produces
$\mathbf{g} = \mathrm{softmax}(W_r \mathbf{h}^{(0)}_s) \in \Delta^{E-1}$, the
token is assigned to $e^\star = \arg\max_e g_e$, and \emph{only} that expert is
evaluated. The assignment is \emph{persisted}: a token keeps one expert for its
whole recursion trajectory, so the block below reuses one expert's parameters at
every depth and one halt head observes a coherent sequence of states. (An earlier
configuration re-drew both per step; it is retained as a labelled ablation and
discussed in \S\ref{sec:honesty}.)

\begin{equation}
\mathbf{h}^{(t+1)} \;=\; \mathrm{LayerNorm}\!\left(
  \mathbf{h}^{(t)} + g_{e^\star}\cdot \mathrm{FFN}_{e^\star}\!\left(\mathbf{h}^{(t)}\right)
\right), \qquad t = 0,\dots,T-1 .
\label{eq:block}
\end{equation}

The gate probability $g_{e^\star}$ multiplies the expert output, which is what
carries task-loss gradient into the router: the discrete $\arg\max$ is not
differentiated, but the scalar scaling the selected expert's contribution is.
There is one block, so depth is computation through time rather than a stack of
depth-specific layers, and no capacity limit is imposed---every routed token is
evaluated by its expert and none is dropped.

\begin{figure}[t]
    \centering
    \includegraphics[width=0.86\linewidth,keepaspectratio]{fig_architecture.png}
    \caption{The MoRE block. A token is routed once to one of $E$ independent
    expert FFNs; that expert's weights are then reused for an adaptive number of
    recursion steps under ACT halting. Setting $T=1$ recovers MoE exactly and
    setting $E=1$ recovers MoR; the input representation is bit-identical in all
    three cases.}
    \label{fig:architecture}
\end{figure}

\subsection{Halting and the training objective}
\label{sec:halting-method}

Halting follows ACT \citep{graves2016act} as used in the Universal Transformer,
with one halt head per expert. Per token, at step $t$:

\begin{align}
p_t &= \sigma\!\left(\mathbf{w}_{e^\star}^{\!\top}\mathbf{h}^{(t)}\right),
  & c_t &= \sum_{\tau \le t} p_\tau ,
  \label{eq:halt-p} \\[2pt]
\text{halt when } &\; c_t + p_t > 1 - \epsilon, & \epsilon &= 0.01,
  \label{eq:halt-cond} \\[2pt]
R &= 1 - c_t \ \text{(remainder)}, &
  w_t &= \begin{cases} p_t & \text{running} \\ R & \text{at the halt step}\end{cases}
  \label{eq:halt-w}
\end{align}

and the token's output is the convex combination
$\mathbf{o}_s = \sum_t w_t \mathbf{h}^{(t)}$, whose weights sum to exactly $1$.
That sum-to-one property is the mechanism by which halting is trained: because
$w_t$ scales how much of each step's state reaches the output, changing $p_t$
changes the task loss, so a differentiable path runs from $\mathcal{L}_{\text{task}}$
into the halt parameters. The boolean comparison in
Eq.~\ref{eq:halt-cond} decides dispatch only. Tokens still active at $t=T$ are
forced out; halted states are frozen and not recomputed, so a halted token costs
nothing at deeper steps.

Depth is trained through the expected-depth term
\begin{equation}
\rho \;=\; \mathbb{E}_{\text{tokens}}\big[\,N + R\,\big],
\label{eq:ponder}
\end{equation}
where the step count $N$ is discrete and contributes no gradient while the
remainder $R$ does. Eq.~\ref{eq:ponder} is the ponder cost. It is worth noting
that the gradient reaching the halt heads through $R$ is small compared with the
task gradient: a direct measurement in our implementation puts the halt gradient
at roughly $285\times$ below the task term. This ratio is the quantitative form
of the failure analysed in \S\ref{sec:mechanism}: the ponder term is large enough
to be optimised away but too weak to hold depth down once the task gradient
accumulates, and once the halting sigmoid saturates the gradient through $R$
vanishes \citep{graves2016act}. \citet{banino2021pondernet} is the principled fix
and is future work (\S\ref{sec:limitations}).

The training objective is a weighted sum of terms that we report separately and
never as a single number:
\begin{equation}
\mathcal{L} \;=\;
\underbrace{\mathcal{L}_{\text{task}}}_{1.0} \;+\;
0.001\,\mathcal{L}_{\text{bal}} \;+\;
0.001\,\rho \;+\;
\underbrace{0.0\cdot\mathcal{L}_{\text{fam}}}_{h\,\text{detached}} \;+\;
\underbrace{0.0\cdot\mathcal{L}_{\text{route}}}_{\text{off}} .
\label{eq:loss}
\end{equation}
$\mathcal{L}_{\text{task}}$ is per-token cross-entropy in nats.
$\mathcal{L}_{\text{bal}} = -H(\bar{\mathbf{g}}) + E\sum_e f_e \bar{g}_e$ combines
a router-entropy bonus with the Switch-style auxiliary \citep{fedus2022switch},
where $f_e$ is the dispatched load fraction and $\bar{\mathbf{g}}$ the mean soft
router probability; it is averaged over the depth calls actually made, so its
magnitude does not grow with $T$. Three terms are structurally zero on language
and the reason is not convenience. $\mathcal{L}_{\text{fam}}$ reads
$\mathbf{h}.\mathrm{detach}()$, so its gradient cannot reach the trunk---which
\emph{strengthens} every specialisation claim below relative to a comparable
synthetic study, where a family cross-entropy is typically one of the larger terms
in the objective. $\mathcal{L}_{\text{route}}$ and the halting-supervision term are
zero because natural language supplies no per-token oracle expert and no per-token
oracle depth to supervise against. The canonical router is therefore unsupervised
with respect to the per-token routing decision.

\subsection{One engine, three modes}
\label{sec:engine}

MoE, MoR and MoRE are configurations of one implementation with an explicit
\texttt{architecture} mode, not three codebases. They share one data pipeline, one
optimiser setup and one evaluation path, so a difference between them cannot be a
difference in the harness. Only the fields their definitions require differ:
$E \in \{6,1,6\}$, $T \in \{1,7,7\}$, MoR's FFN multiplier (below), and the two
loss weights for quantities that do not exist in a given arm (no ponder cost at
depth 1, nothing to balance with one expert).

We write \textbf{MoR} for our recursion-only configuration---one shared block with
adaptive halting---which is in the spirit of weight-tied recursive transformers
with learned halting \citep{dehghani2019ut,bae2025mor} but is not a reproduction
of any specific published system; all three arms are configurations of the same
engine.

\paragraph{Matching the budget.} MoR has one FFN where MoRE has six, so matching
the per-FFN multiplier would compare 5.6M parameters against 0.9M and call the
difference architectural. We instead equalise \emph{total} FFN hidden width,
giving MoR's single FFN a multiplier of $24 = 4 \times 6$. The result is
5{,}584{,}908 parameters for MoE and MoRE, identical, and 5{,}581{,}063 for MoR---a
0.069\% deficit that is six per-expert halt heads, the router and FFN biases. That
residual is irreducible at any integer multiplier, and it runs \emph{against} the
headline result: the arm with slightly fewer parameters is the one that wins in
domain.
